\section{Power Set: Print All Subsequences}
\label{sec:rec-power-set}

\noindent We now begin the chapter's second family: \textbf{exploration
recursion}, or backtracking. This first problem introduces the single
template -- the \textbf{pick/not-pick} decision tree -- that every other
problem in this group (Sections~\ref{sec:rec-power-set}--\ref{sec:rec-letter-combinations})
specializes or extends.

\problemstatement{Given an array (or string) of $n$ distinct elements,
generate \textbf{every} subsequence (equivalently, every subset, since
order within a subsequence here does not matter for counting purposes)
-- all $2^n$ of them, including the empty subsequence.}

\begin{patternbox}
\emph{``Generate all subsequences/subsets''} is the foundational
backtracking pattern: at every element of the input, there are exactly
\textbf{two} choices -- include it in the current subsequence, or
don't -- and these two choices are completely independent across
different elements. A recursion that branches into exactly these two
calls at every position, one per element, visits every one of the
$2^n$ combinations exactly once.
\end{patternbox}

\subsection*{Understanding the problem carefully}

Think of building a subsequence by walking through the input array left
to right, making one binary decision per element: ``is this element in
my subsequence, or not?'' After making this decision for every one of
the $n$ elements, exactly one specific subsequence has been fully
determined. Since there are $2$ choices at each of $n$ independent
positions, there are exactly $2^n$ distinct sequences of choices --
matching exactly the $2^n$ subsequences a set of $n$ elements has.

\begin{ideabox}[The Pick/Not-Pick Template]
$\textsc{Solve}(\textit{index}, \textit{current})$: if
$\textit{index} = n$ (every element has been decided): record
$\textit{current}$ as one complete subsequence; return (base case).
Otherwise: \textbf{pick} -- append $\textit{arr}[\textit{index}]$ to
$\textit{current}$, recurse with $\textsc{Solve}(\textit{index}+1,
\textit{current})$, then \textbf{undo} the append (remove the last
element, restoring $\textit{current}$ to how it was before this branch).
\textbf{Not-pick} -- recurse with $\textsc{Solve}(\textit{index}+1,
\textit{current})$ unchanged (the element is simply skipped).
\end{ideabox}

\subsection*{Why the explicit ``undo'' step after the pick branch is essential}

If $\textit{current}$ is a single shared mutable list passed by
reference into both branches (rather than a fresh copy made for each),
then after the pick branch's entire recursive exploration finishes and
control returns to this level, $\textit{current}$ must be restored to
exactly the state it had \emph{before} the pick was made -- otherwise the
upcoming not-pick branch would incorrectly start from a list that still
contains the picked element. This pattern -- \textbf{choose}, recurse,
\textbf{un-choose} -- is the defining rhythm of every backtracking
algorithm in this chapter; forgetting the un-choose step is the single
most common bug when implementing backtracking with a shared, mutated
data structure instead of fresh copies.

\subsection*{Algorithm}

\begin{enumerate}
  \item $\textsc{Solve}(\textit{index}, \textit{current})$:
  \begin{enumerate}
    \item If $\textit{index} = n$: add a copy of $\textit{current}$ to
          the result list; return.
    \item Append $\textit{arr}[\textit{index}]$ to $\textit{current}$.
    \item $\textsc{Solve}(\textit{index}+1, \textit{current})$.
    \item Remove the last element of $\textit{current}$ (undo).
    \item $\textsc{Solve}(\textit{index}+1, \textit{current})$
          (not-pick branch, $\textit{current}$ unchanged from before
          step 2).
  \end{enumerate}
  \item Call $\textsc{Solve}(0, [\,])$ to start.
\end{enumerate}

\subsection*{Dry run}

\dryrun{Take $\textit{arr} = [1,2]$ ($n=2$).}

The recursion tree:
\begin{itemize}
  \item $\textsc{Solve}(0, [])$:
  \begin{itemize}
    \item Pick $1$: $\textsc{Solve}(1, [1])$:
    \begin{itemize}
      \item Pick $2$: $\textsc{Solve}(2, [1,2])$: base case, record
            $[1,2]$.
      \item Undo; Not-pick $2$: $\textsc{Solve}(2, [1])$: base case,
            record $[1]$.
    \end{itemize}
    \item Undo; Not-pick $1$: $\textsc{Solve}(1, [])$:
    \begin{itemize}
      \item Pick $2$: $\textsc{Solve}(2, [2])$: base case, record
            $[2]$.
      \item Undo; Not-pick $2$: $\textsc{Solve}(2, [])$: base case,
            record $[]$.
    \end{itemize}
  \end{itemize}
\end{itemize}

Recorded subsequences, in the order generated: $[1,2], [1], [2], []$ --
all $2^2=4$ subsequences of $\{1,2\}$, each exactly once.

\subsection*{C++ implementation}

\begin{lstlisting}
#include <bits/stdc++.h>
using namespace std;

void solve(int index, vector<int>& arr, vector<int>& current,
           vector<vector<int>>& result) {
    if (index == (int)arr.size()) {
        result.push_back(current);
        return;
    }

    current.push_back(arr[index]);      // pick
    solve(index + 1, arr, current, result);
    current.pop_back();                 // undo

    solve(index + 1, arr, current, result); // not-pick
}

vector<vector<int>> powerSet(vector<int>& arr) {
    vector<vector<int>> result;
    vector<int> current;
    solve(0, arr, current, result);
    return result;
}

int main() {
    vector<int> arr = {1, 2};
    for (auto& subset : powerSet(arr)) {
        cout << "[ ";
        for (int x : subset) cout << x << " ";
        cout << "] ";
    }
    cout << endl;
    return 0;
}
\end{lstlisting}

\begin{complexitybox}
\textbf{Time Complexity.} There are $2^n$ leaves in the recursion tree
(one per subsequence), and building each subsequence costs up to $O(n)$
to copy it into the result, giving
\[
  O(n \cdot 2^n).
\]
\textbf{Space Complexity.} $O(n)$ for the recursion depth and the
$\textit{current}$ list at any point in time (excluding the output
storage, which is $O(n \cdot 2^n)$ to hold every generated subsequence).
\end{complexitybox}

\begin{takeawaybox}
The pick/not-pick recursion tree, with an explicit choose/recurse/un-choose
rhythm, is the master template for this entire group of problems.
Combination Sum, Subset Sum, and several other sections ahead are all
this exact same tree, either pruned early (skip branches that can no
longer succeed) or with the choice set widened beyond a simple binary
pick/skip.
\end{takeawaybox}
